% Final project report, course 52056 Data Science Project.
%
% Build:
%   cd empathy-classifier
%   python docs/make_figures.py                          % writes docs/_paper_figs/*.png
%   pdflatex -output-directory=docs docs/final_paper.tex % run twice for references
%
% Overleaf: upload this file together with the _paper_figs/ folder.
\documentclass[11pt,a4paper]{article}

\usepackage[utf8]{inputenc}
\usepackage[T1]{fontenc}
\usepackage{lmodern}
\usepackage[margin=2.5cm]{geometry}
\usepackage{graphicx}
\usepackage{booktabs}
\usepackage{caption}
\usepackage{amsmath}
\usepackage{microtype}
\usepackage{url}
\usepackage[hidelinks]{hyperref}

\captionsetup{font=small,labelfont=bf,skip=6pt}

% Tighter section spacing without pulling in titlesec, so this compiles on a
% minimal TeX install as well as on Overleaf.
\makeatletter
\renewcommand\section{\@startsection{section}{1}{\z@}%
  {-12pt \@plus -3pt \@minus -2pt}{6pt}{\normalfont\large\bfseries}}
\renewcommand\subsection{\@startsection{subsection}{2}{\z@}%
  {-9pt \@plus -2pt \@minus -1pt}{4pt}{\normalfont\normalsize\bfseries}}
\makeatother
\setlength{\emergencystretch}{1.2em}  % lets TeX stretch a stubborn line instead of overflowing
\setlength{\parskip}{5pt}
\setlength{\parindent}{0pt}
\graphicspath{{_paper_figs/}{docs/_paper_figs/}}

\title{\vspace{-1.2cm}\Large Predicting the Dominant Empathy Type of a Written Response}
\author{\normalsize Shaul Tolkowsky, Guy Shiff \\[2pt] \normalsize 52056 Data Science Project
\\[2pt] \small Code: \url{https://github.com/shaulTol/LLMES}}
\date{}

\begin{document}
\maketitle
\vspace{-1.0cm}

\section{Abstract}

Rubin et al. \cite{rubin2025} work with a three-way split of empathy: a cognitive
component, understanding what another person is going through; an affective component,
feeling with them; and a motivational component, caring enough to want to help. We take
that split and train a classifier to predict which of the three dominates a written reply
to someone's personal disclosure.

Their study asked whether empathy is valued differently depending on whether people believe
it came from a person or from an AI. Participants read short disclosures together with the
replies written to them, and rated every reply on all three components, five questionnaire
items per component. Each reply therefore carries one participant's scores, and normalising
the three component means leaves a share of cognitive, affective and motivational empathy
that sums to one. The largest of the three is what we call the dominant type. Those shares turn out close to even, with the dominant component
usually holding a little more than a third of the mass, so the task is harder than a
three-way choice sounds. We train on Studies 1 and 1b, which give 1218 replies, and test on
Study 3, which gives 1172 replies written for different prompts. Cognitive dominates 73\% of
the test replies, so always answering Cognitive already scores 0.73 accuracy, and we report
macro F1 throughout with accuracy kept only as a sanity check.

We compare two models on two feature sets. The Semester A baseline freezes a DistilBERT
encoder and trains a small head on its \texttt{[CLS]} vector. The Semester B final model
adds LoRA adapters so the encoder itself can adapt. Each is run on the reply alone and on
the reply concatenated with the disclosure it answers. The final model improves on the
baseline under both feature sets, from 0.336 to 0.400 macro F1, against 0.282 for the
majority-class floor. The two changes interact, since the
disclosure helps only once the encoder can adapt to it. The gains stay small in absolute
terms, and our analysis puts the reason in the data rather than in the model. Applying
ProxySPEX \cite{proxyspex} recovers the same three-word opener features from two different
encoders, and every intervention we tried that removes those features costs macro F1 instead
of forcing the model onto better ones. Swapping DistilBERT for RoBERTa-base lands in the same
range, and on the Affective replies, where the ratings are closest to a three-way tie, all
three of our models fail together. Taken together this points towards the labels rather than
model capacity as what limits the result.

\section{Exploratory analysis (Semester A)}

Every reply in the dataset comes paired with the disclosure it answers, which we call the
\emph{story}, and with three numbers describing how one participant scored it on the
cognitive, affective and motivational components. Those numbers reach us already aggregated:
each participant answered a 15-item questionnaire, five items per component, and the
supplementary data reports the per-component means rather than the individual items. We take
them as given and divide each by the sum of the three, which turns a triple such as
$[9.0, 8.6, 8.6]$ into shares $[0.34, 0.33, 0.33]$. This is what makes the label a
probability distribution and therefore a valid target for the soft cross-entropy of
Section~\ref{sec:train}. It also discards the overall intensity of the ratings, so a reply
scored $[9.0, 8.6, 8.6]$ and one scored $[3.0, 2.9, 2.9]$ become almost the same target and
the model learns the balance between the three components rather than how empathic the reply
was judged to be. That suits the question we ask, which is which component dominates. Two
features of the data shaped every modelling decision that followed, and
Figure~\ref{fig:eda} shows both.

\begin{figure}[h]
\centering
\includegraphics[width=\textwidth]{fig1_eda}
\caption{Two properties of the labels. Panel (a): the dominant type usually holds only a
little more than a third of the rating mass, so most replies are judged a mixture instead
of a clear single type. Panel (b): Cognitive dominates in both splits and more so at test,
which rewards a model that simply learns the training prior.}
\label{fig:eda}
\end{figure}

\textbf{The dominant type barely dominates.} A typical label looks like
$[0.36, 0.31, 0.33]$. Participants rarely scored a reply high on one component alone, and we read this as real
rather than as noise. A reply saying ``that sounds exhausting, and I am
here if you need anything'' genuinely carries perspective-taking and an offer of help at
the same time. It also limits what any classifier can reach, because picking the largest
of three near-equal numbers is close to arbitrary.

\textbf{The class prior is skewed, and it moves between splits.} Cognitive leads in both,
at 67\% of training replies and 73\% of test replies. This drove our choice of metric.
Always answering Cognitive scores 0.734 accuracy and 0.282 macro F1, so accuracy rewards a
model for leaning on the prior, while macro F1 exposes it.

\textbf{The text is short and templated.} Replies run 45 to 110 words and were produced
under a handful of prompt conditions. Word-frequency and style analysis showed the same
few openings reused constantly, among them ``It sounds like\ldots'' and ``I'm really sorry
to hear\ldots'', with lexical diversity low and roughly equal across the three types. We
noted at the time that a classifier might key on those openings. Section~\ref{sec:analysis}
shows that it did.

\section{Training and evaluation process (Semester A)}
\label{sec:train}

\textbf{Splits.} Training uses Studies 1 and 1b less 50 held-out replies per study, giving
% TODO(test-protocol): state here that validation was used only for early stopping,
% and that model selection ran against Study 3.
1218 rows. Those 100 held-out rows form the validation set. Study 3 is the test set, all
1172 rows. We tested on a separate study instead of a random split because we care whether
the model handles new prompts, and a random split would have measured something easier.
The reported numbers are lower as a result.

\textbf{Objective.} Soft cross-entropy against the full rating distribution. We keep the
distribution intact because collapsing it to a single label throws away information about
how close the three types were, and Section~\ref{sec:learned} shows what that costs.

\textbf{Metric.} Macro F1 over the three classes, computed on the argmax of the predicted
distribution. We chose macro over micro averaging because Affective and Motivational cover
only 11\% and 15\% of the test set, and micro averaging would let good Cognitive
performance mask complete failure on the other two. Accuracy appears alongside it only as
a check against the majority-class floor.

\textbf{Protocol.} Early stopping on validation loss with patience 15 and a cap of 100
% TODO(test-protocol): "a single evaluation" is false for the project as a whole.
% Selection across the whole search used Study 3 macro F1; validation was early stopping only.
epochs, then a single evaluation on Study 3. Every
cell of Section~\ref{sec:models} is one training run at a fixed seed, and the confusion
matrix we show for that cell comes from the same run, so the table and the figure describe
one object rather than two. Seed 9 is the seed our Semester B results were presented on and
we keep it throughout, with one exception noted in Appendix~\ref{app:prov}. Runs on this
dataset do vary by seed, so gaps of a few thousandths between cells should not be read as
real; the differences we draw conclusions from are an order of magnitude larger than that.

\textbf{Reference points.} Table~\ref{tab:refs} collects the floors that make the rest of
the report readable. Without them a model at 0.68 accuracy looks like it works when it is
sitting below the trivial floor.

\begin{table}[h]
\centering
\small
\caption{Reference points on the Study 3 test set, ordered by macro F1. The linear baseline
clears the permutation null by 0.027 and a hand-written opener regex by less than that.}
\label{tab:refs}
\begin{tabular}{lcc}
\toprule
Reference & Accuracy & Macro F1 \\
\midrule
Always answer Cognitive & 0.734 & 0.282 \\
Permuted training labels & 0.611 & 0.309 \\
Hand-written regex over opening phrases & 0.567 & 0.343 \\
Linear baseline: frozen \texttt{[CLS]} with linear head & 0.684 & 0.336 \\
\bottomrule
\end{tabular}
\end{table}

The permutation null is the most useful of these. Training the same architecture on
shuffled labels still reaches 0.309, so the linear baseline's 0.336 contains only about
0.027 of learned signal. The regex row sharpens this further: a dozen hand-written patterns
over opening phrases reach 0.343, slightly above the linear baseline, with no learned
representation involved at all. Note also that accuracy and macro F1 disagree across the
whole table. The highest accuracy belongs to the row that answers Cognitive every time and
scores the worst macro F1 of any, which is the clearest argument for the metric we chose.

We call the bottom row of Table~\ref{tab:refs} the \emph{linear baseline} throughout. The
frozen encoder with an MLP head introduced next is a stronger and separate model, and we
call it the \emph{frozen+MLP baseline}.

\section{Prediction models}
\label{sec:models}

\textbf{Baseline model (Semester A).} A frozen DistilBERT encoder produces one
\texttt{[CLS]} vector per text, and a trained linear head turns it into a distribution over
the three types. The encoder never updates, and the baseline uses neither balanced sampling
nor augmentation, so it is the plain configuration; those two are the first steps of the
search in Appendix~\ref{app:search} rather than part of the baseline. This is the model our
Semester A work settled on and the one we carried into Semester B as the comparison point,
and it is the linear baseline of Table~\ref{tab:refs}. We also built a stronger frozen variant with an MLP-256
head, and we return to it below because it matters for how much credit fine-tuning deserves.

\textbf{Final model (Semester B).} The same encoder carrying LoRA adapters of rank 4 on the
query and value projections of all six attention blocks, about 74k trainable parameters
against 66M frozen. A skip connection gives the head both the adapted and the original
\texttt{[CLS]} vector. Full configuration: AdamW, head learning rate 3e-4 and adapter
learning rate 3e-5, weight decay 0.01, batch size 32, MLP-256 head at dropout 0.5, balanced
sampling, and at most 100 epochs with early-stopping patience 15, which in practice stops
around epoch 30. Balanced sampling here means a weighted sampler that draws training rows
with replacement, weighting each row inversely to how common its dominant type is. Batches
therefore come out roughly even across the three types rather than following the 67/12/21
split of the training data. A greedy search that changed one component at a time over roughly 30
steps produced this configuration, and Appendix~\ref{app:search} summarises it.

\textbf{Feature sets.} \emph{Response only} is the 768-di\-men\-sion\-al \texttt{[CLS]}
vector of the reply on its own. \emph{Story+Response} encodes the disclosure and the reply
separately and concatenates them into a 1536-dimensional vector, so the head can see the
situation the reply was written for. Each was tuned on its own learning rate, since the two
differ in input dimensionality. We tried a third option as well, replacing DistilBERT with
RoBERTa-base so the embedding itself changed rather than its input, and it reached
0.376. That sits inside the same band as everything else, so we
report it as a negative result in Appendix~\ref{app:search} and keep the two
\texttt{[CLS]} feature sets in the grid.

Table~\ref{tab:grid} gives all four combinations under macro F1.

\begin{table}[h]
\centering
\footnotesize
\caption{Macro F1 on Study 3 for two models and two feature sets. The two upper rows are
models, one per row, evaluated under each feature set; every one of those four cells is a
single training run, and the confusion matrix for that same run appears in
Figure~\ref{fig:grid}. The final row and the final column are not models or feature sets but
differences between the cells beside them, which is why they sit outside the rules. The
final model improves on the baseline under both feature sets, and the gain is far larger
with Story+Response. Provenance in Appendix~\ref{app:prov}.}
\label{tab:grid}
\begin{tabular}{@{}lcc@{\qquad}c@{}}
\toprule
& \multicolumn{2}{c}{\textit{Feature set}} & \textit{Difference} \\
\cmidrule(lr){2-3} \cmidrule(l){4-4}
Model & Response only & Story+Response & Gain from story \\
\midrule
Baseline: frozen \texttt{[CLS]}, linear head & 0.336 & 0.333 & $-0.003$ \\
Final: LoRA r4 qv all-6 & 0.352 & \textbf{0.400} & $+0.048$ \\
\midrule
\textit{Difference:} gain from LoRA & $+0.016$ & $+0.067$ & \\
\bottomrule
\end{tabular}
\end{table}

\textbf{Reading the grid.} The final model improves on the baseline under both feature sets,
by 0.016 on the reply alone and by 0.067 once the disclosure is included, reaching 0.400
against the baseline's 0.336. The two axes interact rather than adding
up. Story context does nothing for the frozen linear head, which loses 0.003 when given it,
and helps the fine-tuned model by 0.048. That pattern makes sense: a linear map over a
1536-dimensional concatenation has twice the parameters and no way to relate the two halves,
while LoRA can adapt the encoder so the story representation carries something the head can
use. The combination of fine-tuning and story context is therefore worth more than either
alone, which is the main modelling result of Semester B.

\textbf{One qualification on how much credit fine-tuning deserves.} The strongest frozen
model we ever built was not the linear baseline but a frozen encoder with an MLP-256 head on
Story+Response features and latent-Gaussian augmentation to 2500 examples per class, which
reaches 0.377. That model gets its class balance from the augmentation rather than from a
sampler, so both of our stronger models correct the skewed prior in some way and neither
result should be read as coming from architecture alone. Measured against it rather than
against the Semester A baseline, the LoRA run is ahead by 0.023 rather than by 0.067. We
report this because it changes the interpretation: fine-tuning does beat the baseline we
carried forward, and a good part of that gain is reachable with a better head on a frozen
encoder. Section~\ref{sec:analysis} explains why,
and the explanation is the same for both models.

\textbf{Which model we prefer.} The LoRA Story+Response model, as the final model, on the
grounds that it has the best macro F1 of the four cells. It costs more to train, since the
encoder no longer runs once into a cache. Where training cost matters more than the last few
hundredths of macro F1, the frozen MLP head on Story+Response features is the practical
choice.

Figure~\ref{fig:grid} shows a confusion matrix for each cell, row-normalised so the
diagonal reads as per-class recall. Each panel is the run reported in the matching cell of
Table~\ref{tab:grid}. The baseline panels make its weakness visible: the linear head answers Cognitive
almost everywhere and recovers close to nothing on Affective, while the LoRA panels spread
their answers across all three classes.

\begin{figure}[h]
\centering
\includegraphics[width=0.95\textwidth]{fig2_grid_cms}
\caption{Row-normalised confusion matrices, so each diagonal entry is a per-class recall,
for the four cells of Table~\ref{tab:grid}. Every cell over-answers Cognitive and recovers
Affective worst, and that is where the macro F1 goes.}
\label{fig:grid}
\end{figure}

\section{Model analysis}
\label{sec:analysis}

We analyse the linear baseline and the final LoRA model by perturbing their inputs and
watching what the perturbation costs.

If the opener drives the answer, damaging the opener should hurt more than damaging
anything else. We tested this two ways during training. The first removes the first ten
words of every reply, and we ran it on both models, so it supports a comparison between them
(Table~\ref{tab:perturb}). The second replaces those words with an opener taken from a reply
of a different class, at three rates, and we ran it on the final model only, so it speaks to
that model alone (Table~\ref{tab:swap}).

\begin{table}[h]
\centering
\small
\caption{Removing the opener, applied at training and evaluation time. Each column is read
against its own control, and the two controls differ; see Appendix~\ref{app:prov}. Stripping
the opener costs both models real F1 and costs the fine-tuned encoder more than twice as
much.}
\label{tab:perturb}
\begin{tabular}{lcc}
\toprule
& Linear baseline & Final (LoRA) \\
\midrule
None (control) & 0.354 & 0.368 \\
Strip the first 10 words & 0.339 ($-0.015$) & 0.332 ($-0.036$) \\
\bottomrule
\end{tabular}
\end{table}

\begin{table}[h]
\centering
\small
\caption{Replacing the opener with one drawn from a reply of a different class, at three
rates, on the final model. Light swapping barely registers; breaking the opener-to-label
mapping outright at $p=1.0$ costs almost as much as deleting the opener altogether.}
\label{tab:swap}
\begin{tabular}{lc}
\toprule
& Final (LoRA) \\
\midrule
None (control) & 0.368 \\
Swap opener across classes, $p=0.3$ & 0.363 ($-0.005$) \\
Swap opener across classes, $p=0.5$ & 0.366 ($-0.002$) \\
Swap opener across classes, $p=1.0$ & 0.340 ($-0.028$) \\
\bottomrule
\end{tabular}
\end{table}

\textbf{Damaging the opener hurts both models.} Stripping it costs the linear baseline
0.015 and the final model 0.036, and breaking the opener-to-label mapping entirely with a
cross-class swap costs the final model 0.028. Set against the 0.027 of learned signal the
linear baseline holds over the permutation null, these are large. The reading the data
support is that the fine-tuned encoder depends on the opener more than the frozen one does.
Measured against the always-Cognitive floor of 0.282, breaking the opener mapping
takes the LoRA model from 0.368 to 0.340, which removes about a third of its above-floor
signal. The opener carries the most weight of any single feature, and the body of the reply
still contributes the rest.

\textbf{The opener may be a shallow feature the models lean on too heavily, but removing it
does not help.}
The results above look like over-reliance on opener templates, so the obvious remedy is to
train the shortcut out. It does not work. Stripping the opener from
every reply costs the fine-tuned model 0.036, and swapping openers across classes costs
0.005 at $p=0.3$ and 0.028 at $p=1.0$. The one opener intervention that helps is the
same-class swap used in the final recipe, worth 0.007, and that one leaves the
template-to-label mapping intact while varying the wording. Read against the
majority-class floor of 0.282, the stripped model still holds 0.050 of the 0.086 the control
holds above the floor, so the body of the reply carries most of what the model knows even
though the opener is the single most load-bearing feature. Removing the shortcut does not
reveal better features underneath; it removes signal and leaves the rest where it was.

\section{Research task: ProxySPEX feature recovery}

\textbf{The method.} ProxySPEX \cite{proxyspex} explains a model by treating its output as
a function over binary masks of the input tokens and extracting the largest Fourier
coefficients of that function. A coefficient on a token subset $T$ measures what those
tokens contribute jointly, so the method reports interactions between tokens instead of
per-token attributions. The exact spectrum is exponential to compute, so ProxySPEX fits a
gradient-boosted-tree proxy on (mask, output) pairs and reads the coefficients off the
proxy in closed form, using about $O(n \log n)$ model queries.

\textbf{Why it fits our problem.} Section~\ref{sec:analysis} points at opening phrases, and
it does so indirectly, by deleting text and measuring what that costs. It does not identify
which tokens the model combines. Opener templates are multi-word units, so a
method built around token interactions is the right instrument here. We also wanted a test
we could apply unchanged to different encoders.

\textbf{How we applied it.} We ran ProxySPEX over the first ten tokens of each Study 3
reply, allowing subsets up to size three and keeping the top ten coefficients per example.
Two models ran at full budget, the LoRA final model and a RoBERTa-base variant trained with
the same recipe, each on 20 examples per class with 256 masks. The linear baseline ran only
as a reduced smoke test on 4 examples per class with 100 masks, and we treat it as
indicative. The implementation is \texttt{src/proxyspex\_opener.py}.

\textbf{Results.} Three-token subsets dominate the recovered spectrum. For the LoRA model's
Cognitive examples the 200 retained coefficients split into 127 triplets, 62 pairs and 11
singletons, and the other classes and the RoBERTa run look the same. Two things follow.
The dominance of triplets is mechanical: an opener such as ``I'm / really / sorry'' is a
three-word unit, so removing any one word breaks the template and the variance of the
answer collects in three-word coefficients, while no single word shifts the output much
alone. More importantly, two architecturally different encoders converge on the same
features, with the reduced baseline run pointing the same way. We would not expect a quirk of
one model to reproduce across encoders, so the more likely reading is that the labels reward
this structure, though two encoders is a thin basis for saying so.

\textbf{Do the models fail on the same replies?} We compared the test errors of the three
reference models against an independence null. On Affective rows all three are wrong
together 107 times out of 132, and independence already predicts 103.4, so this reflects
each model having a marginal error rate above 83\% on that class. On Cognitive rows all
three are wrong on only 11 of 860 against a null of 1.0, $p \approx 10^{-8}$, which is
genuine shared confusion over a small set of ambiguous rows. On Motivational rows all three
are wrong 111 times out of 180 against a null of 96.0, $p = 0.015$, so shared confusion
appears there too, spread across a large fraction of the class. For Affective the labels
themselves look like the likeliest explanation: mean label entropy on Affective rows is 1.57 bits of a
possible 1.585, so those rows are near-uniform mixtures with no recoverable answer.

\section{Conclusions}

\textbf{What worked.} Treating macro F1 and the permutation null as the real metrics kept
us honest. The linear baseline sits at or below the 0.734 trivial floor on accuracy, and only
the macro F1 view showed it. Optimising one metric is fine, but reading only that metric
hides what the model is doing: our baseline checkpoint scores 0.719 accuracy and 0.327 macro
F1, and the two even move in opposite directions when we add Gaussian noise to all 768
encoder features, which raises macro F1 to 0.349 while dropping accuracy to 0.611 because the
noise lands on the Cognitive bias rather than on signal. Giving the model the disclosure
alongside the reply helped, though only once the encoder could adapt to it, which took us a
while to see. Correcting the skewed class prior mattered too, though less in the headline
number than in what the model predicts. The sampler barely moves macro F1 on its own, yet
both models that beat the linear baseline correct the prior in some way, one with a sampler
and one with augmentation, and the confusion matrices show the effect plainly as answers
spreading off Cognitive and onto the two minority types. Tuning also paid, more than we expected for something we had treated as
bookkeeping: retuning the learning rate for the Story+Response features moved them from
0.373 to 0.381, and fixing the early-stopping schedule moved the MLP head off 0.330.
Re-running a promising configuration before believing it stopped us several times from
claiming gains that did not hold up. Caching the frozen encoder's outputs turned
a five-minute experiment into a one-second one, and a search of this size would have been
impossible without it.

\textbf{What did not work.} Almost everything on the model side. LoRA, partial encoder
unfreezing, attention and mean pooling, deeper and wider heads, class-weighted and focal
losses, label sharpening, random forests, chunked features and a switch to RoBERTa-base all
landed in the same 0.35 to 0.40 band. Training on hard labels actively hurt, dropping macro
F1 to 0.166. The most useful result of the project is a negative one, and establishing it
with any confidence took about 30 steps of search.

\textbf{The search itself was harder to read than we expected.} A step that helps at one
stage can stop helping at another for reasons that are not obvious in advance.
Story+Response is the clearest case: at the learning rate tuned for response-only features
it scored 0.373 against response-only's 0.374, so we briefly held evidence that the
disclosure was useless, and retuning brought it to 0.381. Our original MLP was stopping
early at a mean of 11.7 epochs, so its 0.330 described the schedule more than the
architecture. Latent augmentation is the case we understand best in hindsight. It helps a
frozen encoder, where the \texttt{[CLS]} vectors are a fixed point cloud and Gaussian noise
genuinely smooths the boundary the head draws through them. Once LoRA is training that cloud
moves at every step, so noise calibrated to the original representation is calibrated to
nothing, and the head cannot separate augmentation noise from the encoder actually learning.
The lever became null and we had recorded it as a gain. Some of our earlier negative results
were therefore not trustworthy at the moment we drew them.

% TODO(test-protocol): add a limitation stating that architecture and hyperparameter
% choices were made against Study 3, so the absolute numbers are optimistic as a measure of
% generalisation and the gaps between cells are more trustworthy than their levels.
\textbf{Limitations.} We never measured the ceiling we keep pointing at, and we could not
have done so with this dataset. By the ceiling we mean the best macro F1 any predictor could
reach given how unstable the target is, and the way to measure it is to have several people
rate the same reply and then score one rater's answer against another's under the same
metric. No model should be expected to predict a label better than a second person does.
Every reply here carries the scores of exactly one participant, so there is no disagreement
to measure. Every statement we make about the ceiling is
therefore an inference from convergent evidence rather than a result, and the resampling
check of Appendix~\ref{app:prov} is a rough indication and nothing stronger. The test set is one study, so generalisation
here means generalisation to a single new prompt set. The numbers we report are single
runs rather than averages, so a few thousandths either way should not be read as meaningful,
and the LoRA response-only cell predates part of the final recipe.

\textbf{Future directions.} The obvious next step is to measure the limit rather than infer
it, and that needs a small collection effort rather than a better model. Having two or three
people label a few hundred replies with the three-way question directly would give a
human-against-human macro F1 and settle where the ceiling sits, and it is cheap next to
re-running the original study. Beyond that, extending the error-overlap analysis to
per-example agreement on the 107 Affective rows all three models miss, and asking whether a
different label schema, multi-label rather than a forced mixture, would carry more
information per reply.

\textbf{What we learned from the process.}\label{sec:learned}
Noisy, limited data sets the ceiling, and no amount of
feature engineering or model sophistication moves it. Swapping a linear head for an MLP, for
LoRA fine-tuning, and finally for a different pretrained encoder moved macro F1 from 0.336 to
0.400, against a majority-class floor of 0.282 that costs nothing at all, and training on
shuffled labels already reaches 0.309. The labels are the reason, and they are hard in two compounding ways. Judging emotional
content in text is difficult for people and not only for models, since whether a reply
understands someone, feels with them or wants to help them are overlapping readings of the
same sentences. On top of that, nobody in this study was ever asked which type a reply
belonged to. Participants answered fifteen questionnaire items, five per component, and the
component scores are averages over those items; we then normalise the three and take the
largest. The target is a construct assembled two steps away from anything a person decided,
each reply is rated by a single participant rather than a panel, and different replies are
rated by different participants, so between-rater variation enters the dataset as noise we
can neither measure nor average away. Affective rows average 1.57 bits of label entropy out
of a possible 1.585, and of 132 Affective test rows exactly one is answered correctly by all
three of our models. Given labels like these, we would not expect any model to do much
better.

\section*{Code availability}
All code for this report, together with one result artefact behind every number in it, is at
\url{https://github.com/shaulTol/LLMES}. The \texttt{main} branch holds the report and the
code that produces it; the \texttt{process} branch holds the full research history,
including every step of the search and the rejected branches summarised in
Appendix~\ref{app:search}.

\begin{thebibliography}{9}
\bibitem{rubin2025} Rubin, M., Li, J. Z., Zimmerman, F., Ong, D. C., Goldenberg, A., \&
Perry, A. (2025). \emph{Comparing the value of perceived human versus AI-generated empathy}.
Nature Human Behaviour. \url{https://doi.org/10.1038/s41562-025-02247-w}
\bibitem{proxyspex} Butler, L. et al. (2025). \emph{ProxySPEX: Inference-Efficient
Interpretability via Sparse Feature Interactions in LLMs}.
\bibitem{lora} Hu, E. J. et al. (2021). \emph{LoRA: Low-Rank Adaptation of Large Language
Models}. ICLR 2022.
\bibitem{roberta} Liu, Y. et al. (2019). \emph{RoBERTa: A Robustly Optimized BERT
Pretraining Approach}.
\bibitem{distilbert} Sanh, V., Debut, L., Chaumond, J., Wolf, T. (2019). \emph{DistilBERT,
a distilled version of BERT}.
\end{thebibliography}

\newpage
\appendix

\section{Architecture search}
\label{app:search}

A greedy search reached the final model, changing exactly one component of the accepted
% TODO(test-protocol): name the metric explicitly as Study 3 macro F1, i.e. test-set selection.
configuration at each step and keeping the change only when macro F1 improved, over about
30 steps. Table~\ref{tab:chain} lists the accepted chain, and
\texttt{docs/architecture\_search.md} on the \texttt{process} branch documents the rejected
branches.

\begin{table}[h]
\centering
\small
\caption{Accepted chain, one change per row, each a single run at seed 9 so the rows compare
with each other and with Table~\ref{tab:grid}. No individual step is large enough to stand
on its own, so read the chain as a direction of travel rather than as a sequence of
separately meaningful gains.}
\label{tab:chain}
\begin{tabular}{@{}lp{7.6cm}c@{}}
\toprule
Configuration & Change from the row above & Macro F1 \\
\midrule
Linear baseline & n/a & 0.336 \\
+ latent augmentation & Gaussian noise on \texttt{[CLS]}, to 2500 per class & 0.335 \\
+ MLP head & linear head becomes MLP-256, GELU, dropout 0.3 & 0.353 \\
+ story features & concatenate story \texttt{[CLS]} with response \texttt{[CLS]} & 0.377 \\
+ LoRA & rank-4 qv adapters, all 6 layers, skip connection & 0.400 \\
\bottomrule
\end{tabular}
\end{table}

\textbf{Rejected directions}, each tested and abandoned: random forests; chunked features
over openers, thirds, quarters and first-$N$ windows; class-weighted soft cross-entropy;
label sharpening; heavy weight decay; deeper and wider MLPs; mean, attention and
cls+mean+max pooling; partial encoder unfreezing of the top one to three layers, none
beating LoRA; L1 regularisation on the LoRA delta, which collapsed to 0.276; and
BERT-fill augmentation, which timed out before completing.

\textbf{A different base encoder.} DeBERTa-v3-base was the intended comparison. It could
not be loaded under transformers 5.x because of a SentencePiece and tiktoken
incompatibility, so we substituted RoBERTa-base, which brings a different pretraining
objective, a BPE tokenizer and roughly ten times the pretraining data of DistilBERT's
teacher. Its best configuration reached 0.376, inside the same
band as everything else.

\section{Additional diagnostics}

\textbf{B1. Opener probe.} Truncating each reply to its first $N$ words, encoding, and
training a linear head gives macro F1 of 0.297, 0.310, 0.308, 0.324 and 0.350 at $N$ of 3,
5, 10, 20 and 50, against 0.354 for the full reply. These are best read against the
majority-class floor of 0.282, since everything below it is free. On that basis the first
three words recover 0.015 of the 0.072 the full reply reaches above the floor, or about a
fifth of the signal, and the first ten recover about a third. Words past 50 add nothing
measurable. An earlier write-up of this probe reported the first three words as reaching
84\% of full-reply F1, which compares the raw scores without subtracting the floor and
overstates the opener considerably.

\textbf{B2. Hard-label training.} Training the final model on one-hot labels taken from the
argmax swings the bias completely. Cognitive recall falls from 79\% to 0\% while Affective
and Motivational rise to 52\% and 67\%, and macro F1 collapses to 0.166. The argmax maps
$[0.40, 0.30, 0.30]$, clearly Cognitive, and $[0.36, 0.34, 0.30]$, a near-tie, onto the same
target, discarding the calibration the ratings carry.

\textbf{B3. Leave-one-out influence.} Influence is spread thin without being negligible.
96\% of training rows have absolute influence above 0.001 macro F1 and only 38\% fall below
0.01, with mean $-0.0034$ and standard deviation 0.0193 across 1218 retrains. A small
contiguous cluster of Study-1 rows is actively harmful, and removing any one of them
improves test macro F1 by roughly 0.04. We did not act on this, because the influence was
ranked against the test set and using it to clean training data would leak. Readers who saw
either presentation should note that both decks quoted the 96\% figure against the wrong
threshold, as ``96\% of train rows have $|$importance$| < 0.01$'', and concluded influence
was negligible. The artefact \texttt{outputs/a4\_summary.md} reports 96.0\% above 0.001,
with only 38.5\% below 0.01. The reading given here is the corrected one.

\section{Further confusion matrices}

\begin{figure}[h]
\centering
\includegraphics[width=0.9\textwidth]{figC1_extra_cms}
\caption{Row-normalised confusion matrices for three further variants. The linear baseline
and the RoBERTa-base variant over-answer Cognitive exactly as every cell of
Figure~\ref{fig:grid} does, despite the encoder change. The hard-label matrix shows most
clearly why the rating distribution matters: the Cognitive row empties completely and macro
F1 falls to 0.166.}
\label{fig:extra}
\end{figure}

\section{Provenance and superseded numbers}
\label{app:prov}

\textbf{A resampling check, reported here rather than in the main text.} Treating each row's
rating shares as a distribution and drawing two independent hard labels from it gives an
agreement of 0.342 macro F1. We ran this hoping for something like a label-noise floor, and
it does not support that reading, for two reasons. Each row holds the scores of one
participant rather than a panel, so the three shares describe how that person split their
ratings across the components and not disagreement between raters; drawing from them
simulates nothing that happened. And our models predict the argmax rather than a draw, which
beats sampling by construction, so a model can exceed 0.342 without contradicting it, and
ours does. We keep the number as a rough indication that the labels sit close to a three-way
tie, and we draw no bound from it. It has no run artefact in \texttt{outputs/}; it was
computed for the Semester B presentation and the script was not retained.

\textbf{Table~\ref{tab:refs}.} The linear-baseline row is the seed-9 run, while the single
checkpoint analysed in Section~\ref{sec:analysis}, at accuracy 0.719 and macro F1 0.327, is
a separate object. The permutation row is the 100-seed mean of
\texttt{src/a1\_permutation\_null.py}, since a null is a distribution by construction. The
regex row comes from \texttt{src/nonneural\_baselines.py}, written for this report because
the original was run for the Semester B presentation and its artefact was not retained. It is
therefore a reconstruction and does not reproduce the presented figure exactly, 0.343 here
against 0.361 then, which is expected since the original pattern list was not kept and ours
is rebuilt from the opener templates the error analysis surfaced. It uses the same split and the same test
rows as every other number in the report.

\textbf{Table~\ref{tab:chain} and the ladder of Section~\ref{sec:learned}.} Every row except
the LoRA one comes from \texttt{src/chain\_single\_seed.py}, recorded in
\texttt{outputs/chain\_single\_seed.json}. The LoRA row is the run behind
Table~\ref{tab:grid}'s best cell. The RoBERTa figure of 0.376 is scored from
\texttt{outputs/preds\_test\_roberta\_winner.npz}.

\textbf{Table~\ref{tab:grid}.} Both baseline cells are seed-9 runs of the linear head on
cached embeddings, recorded in \texttt{outputs/linear\_head\_feature\_sets\_seed9.json}. The
LoRA Story+Response cell is seed 9 of the deployable checkpoint. The LoRA response-only cell
is the one exception to seed 9: that configuration was run at seed 0 and we did not repeat
it at seed 9. It also predates the skip connection and the decoupled learning rates, so it
understates that cell and the $+0.016$ gain beneath it is a lower bound on what LoRA
achieves there. The
frozen+MLP figure of 0.377 quoted in Section~\ref{sec:models} is a seed-9 run of the same
head and cache used by \texttt{src/run\_story\_100seeds.py}.

\textbf{Figure~\ref{fig:grid}.} Every panel is the same run as the matching cell of
Table~\ref{tab:grid}, so the two report one object rather than two. The LoRA Story+Response
panel is the deployable seed-9 checkpoint, which is the run behind that cell's 0.400.

\textbf{Tables~\ref{tab:perturb} and~\ref{tab:swap}.} Unlike Table~\ref{tab:grid}, these are
averages over 10 runs per cell. The perturbation effects are a few hundredths of macro F1, which a single run
cannot separate from run-to-run variation, so averaging is the only way to read them. Each
column is read against its own control, and the two controls differ: neither matches the
single-run figures of Table~\ref{tab:grid}, which is why the columns are read only against
themselves. The swap rows compare within the final model only.

\end{document}
